\section{The joint small-$a$, small-$e$ corner: a bounded-ratio closure}
\label{sec:r14-r9-bounded-ratio}

\begin{theorem}
Retain the full physical and stationary source domain R9 of (WF3),
including all inherited exclusions, $n\ge0$, and independent source
and testing fields. For every local testing pair $u,v\ne0$ for which
$w=v/u$ is bounded in the testing field, the entire original C2 parent
is finite tensor SOS, with every state component retained. Literal
axes are also covered. The theorem does not cover the remaining
local chart $v/u$ unbounded, and therefore is not an all-tests
admission of R9.
\end{theorem}

The proof uses the exact identities (WC3)--(WC17a), not the unit-$a$
contact-lifting argument. Take $u,v>0$ using the inherited exact sign
gauges. The source quantities $B,a,h,e$ are positive infinitesimals,
$t$ is a unit with $t_0>1$, and $r=1-e$. Put
\[
 c=\sqrt{a/t},\quad k=\sqrt{aB},\quad
 A_1=uv,
 \quad A_2=ck\frac{v^2}{u},\quad A_3=\frac{k^2}{c},
 \quad G=kv,\quad F_v=v^2-tB.
\]
All divisions and radicals are pure-field operations. In particular
$w=v/u$ is bounded and $c$ is infinitesimal. The crucial identity is
\[
 \boxed{G^2=(cw)A_1A_3.}                              \tag{R9B1}
\]
Thus the negative circuit term is an infinitesimal off-diagonal
coefficient in the $(A_1,A_3)$ frame. No circuit conductor and no
comparison of source and testing rates is needed.

\paragraph{Joint source regularity at $a=0$.}
Use the unchanged definitions $q_s,\widetilde m,c_F,f_A,f_B,f_C,f_G$
from (WC1), (WC6), and (WC8), and put $f_{C0}=f_C|_{U=0}$.
Unlike $\partial_Uf_C$, all coefficients needed below are jointly
regular at $B=a=h=0,r=1$. More explicitly,
\[
\begin{array}{c|c}
\text{expression}&\text{joint residue}\\\hline
c_F&1/(2t_0)\\
f_A|_{U=V=0},\ f_B,\ f_{C0}&1/2\\
f_G&3/2\\
\widetilde m/B&0\\
K_X/B,\ K_E/B,\ K_Y/B&t_0/2,\ -1/2,\ 0\\
d_4,\ b_4&t_0/[2(t_0-1)],\ 0\\
\alpha/B^2&t_0/2
\end{array}                                                   \tag{R9B2}
\]
These are finite rational identities. To verify them, use the
unit common denominator
\[
 \mathcal D=2rg(1-h)(1-hB)(t-1+h+aBtr)\Delta
\]
from (WC3), insert the exact source numerators, and divide numerator
and denominator by each explicit monic factor $B$ and $a$.
The required numerator remainders vanish identically, not just at
the corner. Each remaining denominator has a nonzero corner value.
The accompanying focused checker performs precisely these divisions
and prints the denominator and residue for every row. There is no
$e/B$, $h/a$, or mixed-rate regularity assumption. All remaining
source coefficients of (WC17a) are rational with the same unit
denominators and are therefore bounded on this chart.

\paragraph{A strict scalar Gram without a contact lift.}
Define the legal separated roots
\[
 T_q=\sqrt{q_s}\,u,\qquad T_n=\sqrt n\,r u^2,
 \qquad T_\alpha=\sqrt\alpha\,u,\qquad
 d=1-\frac a t w^2.
\]
Here $q_s>0$ and $n\ge0$; literal $n=0$ gives the zero root
$T_n$ and requires no division. Also $d$ has residue one.
The exact scalar identity (WC9) becomes
\[
 \mathcal S=d^2(T_q^2+T_n^2)+c_FF_v^2
       +f_AA_1^2+f_BA_2^2+f_{C0}A_3^2
       -f_G(cw)A_1A_3+T_\alpha^2.                 \tag{R9B3}
\]
Indeed $H_u=-du$ and the formerly problematic term is exactly
\[
  (f_C-f_{C0})A_3^2=\alpha u^2.
\]
It is retained as a square, not treated as a bounded perturbation.
Since $f_A=f_A|_{U=V=0}+d_4U+b_4V/t^2$, (R9B2) proves that
(R9B3) has a bounded Gram with strictly positive diagonal real
residue on
\[
 T_q,T_n,F_v,A_1,A_2,A_3,T_\alpha.
\]
Its only displayed off-diagonal entry is infinitesimal by (R9B1).

\paragraph{Every original state normal.}
Use exactly the full-state frame (WC4),
\[
 E=tuW-v,\qquad N_y=(Y-huW)/\sqrt h,
 \qquad \xi=X+ru,
\]
initially at $Z=1$. The original hard block has the unchanged
strict residue (WC5), because its rational coefficients have no
negative power of $a$ and (WF6)--(WF7) supply its source limits.
The extracted angular contributions in (WC15) are infinitesimal
linear combinations of $T_q,T_n$: simply replace
$R_q=-dT_q$ and $R_n=-dT_n$. Their coefficients remain bounded
infinitesimals even when $n=0$.

For the complete residual normals, use their exact polynomial
arrays (WC17a). Every term is either $Buw^j$ with $0\le j\le3$
or $u^iw^j$ with $i\ge3,j\ge1$, with bounded pure-source
coefficients by (R9B2). The following exact allocations have only
infinitesimal coefficients:
\[
 \boxed{
 Buw^j=\frac{u w^{j+1}}t A_1-\frac{u w^j}t F_v,
 \qquad
 u^iw^j=u^{i-2}w^{j-1}A_1.}                         \tag{R9B4}
\]
The first identity follows from $tB=v^2-F_v$ and $A_1=uv$;
the second is immediate. These identities include the formerly
problematic $Bu$ term, and never divide by $a$, $k$, or a mixed
element. Hence all three full first normals are infinitesimal
linear combinations of roots already carrying strict scalar
pivots. The $\sqrt h$ multiplying the $N_y$ normal is bounded
and causes no difficulty. Joining them with the hard block gives
an exact finite bounded Gram with strictly positive real residue.
The inherited bounded-Gram lemma now proves tensor SOS for the
complete parent.

Homogenize by $Z^2$ exactly as after (WC28), so all roots become
homogeneous state-linear expressions and no $Z$ inverse remains.
The resulting identity includes $Z=0$. The original sign gauges
handle all signs, the two exact positive-axis identities (WC28)
handle literal axes, and real/imaginary splitting handles complex
states. No norm, original packet, or hard normal was discarded.

\paragraph{Exact remaining scope.}
Together with the R13 remote reduction, this leaves on R9 only
$u,v\ne0$ both infinitesimal with $\omega=u/v$ infinitesimal.
The new identity (R9B1) does not claim $cw$ bounded when $w$ is
unbounded. That reciprocal local chart remains an explicit proof
obligation. The accompanying 23 exact checks (including twelve
joint source regularity checks) verify the new finite identities;
they are not a replay of an older suite.
